% ID: SEL-D0077-OA-E03
% ESTADO: cotejada. Segunda pasada visual y borrador compilado revisados.
% Fuente: Selectividad/sel_2008_matematicas_aplicadas/Mat CCSS 5 - Examen y criterios de correción.pdf
% SHA-256: 6b6f35fe4872415542ee008d454a2502926c6ed0fffe87639ed13bed50759c6b
% Página: 1; opción A; ejercicio original 3.

\noindent\underline{Parte I}

En un aula de informática hay 20 puestos de ordenador. De ellos, 10 son compartidos y otros 10 son individuales. De los puestos compartidos, hay 3 en los que el ordenador no funciona, de los individuales hay 2 en los que el ordenador no funciona.
\begin{enumerate}[label=\alph*)]
\item \textbf{(1 punto)} Seleccionado al azar un puesto en el aula, \textquestiondown{}cuál es la probabilidad de que no funcione el ordenador?
\item \textbf{(1 punto)} Si se elige al azar un puesto en el que funciona el ordenador, \textquestiondown{}cuál es la probabilidad de que sea compartido?
\end{enumerate}
\par\noindent\begin{minipage}{\linewidth}
\noindent\underline{Parte II}

El peso, en kg, de los alumnos de primaria de un colegio sigue una distribución Normal de media 28 kg y desviación típica $2.7$ kg.

Consideremos muestras aleatorias de 9 alumnos.
\begin{enumerate}[label=\alph*)]
\item \textbf{(0.5 puntos)} \textquestiondown{}Qué distribución sigue la media de las muestras?
\item \textbf{(1.5 puntos)} Si elegimos, al azar, una de esas muestras, \textquestiondown{}cuál es la probabilidad de que su media esté comprendida entre 26 y 29 kg?
\end{enumerate}
\end{minipage}\par
